% !TEX root = ../../draft_v2.tex
%==============================================================================
%  Appendix D4 : Bargaining interpretation of the control premia (m0,m1)
%  ----------------------------------------------------------------------------
%  Drop-in appendix for "Liquidity, Activism Disclosure, and Takeover Premia".
%  Uses the body macros \E, \PP, \1, \Var and biblatex \citet/\citep.
%
%  HONEST RELABEL (this revision).  The body treats the bid-conditional control
%  premia (m0,m1) and the engagement-success probability rho as reduced-form
%  primitives, with tilde m = m0 + rho(m1-m0) and standing Assumption A3 stating
%  tilde m > m0.  Five earlier attempts tried to DERIVE the load-bearing
%  threat-point shift d(1)-d(0)=lambda*Dtil_eng from a fully specified
%  outside-option game so that A3 would follow with NO residual structural
%  assumption.  That program did NOT close: the pivotality/appropriability
%  coefficient lambda -- which carries the entire sign and magnitude of the
%  wedge -- cannot be pinned without a complete tender-game equilibrium that is
%  orthogonal to this paper's disclosure mechanism.  Asserting "lambda>0 hence
%  d(1)-d(0)>0" as a Lemma would merely relabel the assumption as a theorem,
%  exactly the circularity the critique flags (Prop 7 "re-assumes its own
%  conclusion").
%
%  We therefore adopt the strategy of Edmans, Goldstein, and Jiang (2015) and
%  related corporate-finance theory: state the result under an EXPLICIT,
%  economically transparent primitive condition (Condition BG below), motivate
%  that condition with a free-rider tender mechanism (Grossman-Hart 1980) but
%  present the motivation HONESTLY as a heuristic, not a closed derivation, and
%  carry the remaining numerical claim (single-peak survival) as a NUMERICALLY
%  VERIFIED regularity citing the Phase-0 solves -- clearly flagged as a
%  calibration-level regularity, not a proof.  What is genuinely proved is kept
%  as proved (the Nash split, the wedge's xi-invariance, the reduced-form
%  nesting, the rho^2 event tree, the identification non-result).  What remains
%  open is stated precisely (Section app:bg-open).
%
%  ---------------------------------------------------------------------------
%  WHAT CHANGED RELATIVE TO THE EARLIER "MICROFOUNDATION" DRAFT:
%   * The threat-point shift is now a STATED structural Condition (Condition BG),
%     NOT a "derived object" Lemma.  The free-rider tender stage is presented as
%     MOTIVATION (a heuristic that makes the sign plausible), explicitly flagged
%     as not constituting a proof of lambda>0 from deeper primitives.
%   * A3 is a THEOREM CONDITIONAL ON Condition BG (honest: the conditional
%     implication is valid; the antecedent is an assumption, transparently so).
%   * The single-peak comparative static is a NUMERICAL REGULARITY at the
%     Phase-0 calibration, conditional on the chord-curvature condition (C**) of
%     Appendix D.1; it is NOT presented as a theorem.
%   * lambda is carried explicitly through identification and reconciliations;
%     the "sign for any lambda>0" claim is restricted to the SIGN, with magnitude
%     carrying lambda; mtilde-m0 = rho^2(1-theta)*lambda*Delta_eng.
%   * rho^2 justified by an explicit event tree (Lemma lem:bg-tree) with two
%     INDEPENDENT Bernoulli draws B1,B2; formal test E[B1 B2]=rho^2 != rho.
%   * Identification stated honestly: only the wedge PRODUCT is pinned; theta is
%     set-identified; the theta comparative static is an if-identified remark.
%   * Price P reconciliation: P is a bidder-to-float transfer orthogonal to the
%     bidder-vs-bloc pie, so P never enters the pie G.
%  ---------------------------------------------------------------------------
%
%  TWO DISTINCT PRIMITIVES IN THE BODY (the key to the numbers; an earlier draft
%  wrongly fused them into one Delta).  The body's baseline table (draft line
%  1292) sets the STANDALONE-improvement parameter Delta = 0.25, and SEPARATELY
%  the premia m0 = 0.10, m1 = 0.30, so the WEDGE m1 - m0 = 0.20 is its own
%  primitive, NOT a function of the standalone Delta.  Hence:
%        standalone channel : tilde Delta = rho*Delta = 0.9*0.25 = 0.225
%        wedge (premia)     : m1 - m0 = 0.20  (independent primitive)
%        folded premium     : tilde m - m0 = rho*(m1-m0) = 0.9*0.20 = 0.18
%  The bargaining reading introduces its OWN engagement value Delta_eng that
%  drives the threat point and hence the WEDGE; it is conceptually the engagement
%  improvement that reaches the sale, and it need NOT numerically equal the
%  body's standalone Delta = 0.25.  We give consistent reconciliations:
%     (R1) theta=1/2, lambda=1 : Delta_eng = (m1-m0)/((1-theta)*lambda*rho)
%                              = 0.20/0.45 = 4/9
%     (R2) Delta_eng = body Delta = 0.25, lambda=1 :
%          theta = 1 - (m1-m0)/(lambda*rho*Delta_eng) = 1 - 0.20/0.225 = 1/9
%     (R3) theta=1/2, lambda=1/2 : Delta_eng = 0.20/(0.5*0.5*0.9) = 8/9
%  All reproduce the body's wedge 0.20 and (via the body's single-rho fold)
%  tilde m - m0 = 0.18.  There is NO divide-by-rho step anywhere.
%
%  PHASE-0 SINGLE TUPLE cited everywhere in this appendix:
%        (grid, kstar, peak, amplitude) = (99-pt over [0.01,0.99], 0.59,
%         0.07756, 0.0107).
%  Near-endpoints Delta^min(1e-6)=0.07020822, Delta^min(1-1e-6)=0.07020753
%  (gap 6.86e-7), endpoint symmetry as a kappa->0,1 LIMIT (Phase-0 fact sheet).
%  The appendix's own GE re-solve (d4_bargaining_check.py [D1]/[D2]) returns
%  kstar=0.60, peak 0.07755, amp 0.0106 on its coarser solver grid -- the same
%  single hump; [D2] confirms the hump survives a state-dependent wedge for
%  lambda_sd in {-0.4,-0.2,0,+0.2,+0.4}.
%
%  COMPILE NOTE.  The draft preamble defines theorem/proposition/lemma/
%  corollary/definition (and \E,\PP,\1,\Var,\R) but NOT remark/condition/numreg;
%  all three are guarded with \@ifundefined below.  Only standard environments
%  are used.
%==============================================================================

% --- defensive macro guards (no-ops if the body already defines these) -------
\providecommand{\E}{\mathbb{E}}
\providecommand{\PP}{\mathbb{P}}
\providecommand{\1}{\mathbf{1}}
\providecommand{\Var}{\operatorname{Var}}
\providecommand{\Dtil}{\widetilde{\Delta}}
\providecommand{\mtil}{\widetilde{m}}
\providecommand{\mR}{m^{R}}

\makeatletter
\@ifundefined{numreg}{%
  \newtheorem{numreg}{Numerical Regularity}%
}{}
\@ifundefined{c@condition}{%
  \newtheorem{condition}{Condition}%
}{}
\@ifundefined{c@remark}{%
  \theoremstyle{remark}\newtheorem{remark}{Remark}\theoremstyle{plain}%
}{}
\makeatother

\section{A bargaining reading of the control premia under an explicit
threat-point condition}
\label{app:bg}

This appendix gives a generalized-Nash-bargaining reading of the bid-conditional
control premia $(m_{0},m_{1})$ used in the body. We are deliberate about the
\emph{epistemic status} of every claim, because the load-bearing object --- the
engagement-induced shift in the seller bloc's disagreement payoff --- cannot be
honestly derived from deeper primitives within the scope of this paper.

\paragraph{Honest scope statement.}
A full microfoundation would derive the threat-point shift $d(1)-d(0)>0$ (and
hence Assumption~\textup{(A3)}, $\mtil>m_{0}$) from primitives, with \emph{no}
residual structural assumption. We do not achieve that, and we say so plainly.
The obstruction is concrete: the sign and magnitude of the shift are governed by
a pivotality/appropriability coefficient $\lambda$ that can only be pinned by a
complete tender-game equilibrium --- a model orthogonal to this paper's
disclosure mechanism. Writing ``$\lambda>0$, therefore $d(1)-d(0)>0$'' as a
lemma would relabel the assumption as a theorem and re-assume the very thing to
be shown. Following the strategy of \citet{EdmansGoldsteinJiang2015} and the
corporate-finance theory tradition of stating results under transparent
primitive conditions, we instead: (i) \emph{state} the threat-point shift as an
explicit, economically interpretable condition (Condition~\ref{cond:bg}); (ii)
\emph{motivate} that condition with a free-rider tender heuristic in the spirit
of \citet{GrossmanHart1980}, flagged honestly as motivation rather than proof;
(iii) prove what genuinely follows from the bargaining protocol \emph{given} the
condition; and (iv) carry the one remaining numerical claim (the single-peak
comparative static) as a \emph{numerically verified regularity} at the Phase-0
calibration, clearly labelled as such.

\paragraph{What this buys, stated without inflation.}
The reduced form posits an opaque premium ordering, $m_{1}>m_{0}$. The bargaining
reading replaces it with a \emph{transparent} structural condition --- engagement
raises the seller bloc's appropriable no-sale continuation value (its bargaining
threat point) by $\lambda\,\Dtil_{\mathrm{eng}}$ --- whose \emph{sign} is made
plausible by a recognisable free-rider mechanism and whose \emph{magnitude} ties
to interpretable bargaining-power and engagement-improvement primitives. That is
real interpretive progress and it sharpens the comparative statics. It is
\emph{not} a free lunch: the coefficient $\lambda$ remains a reduced-form
primitive (Remark~\ref{rem:bg-threat-honest}), so A3 is a theorem
\emph{conditional on} Condition~\ref{cond:bg}, not an unconditional theorem.

\paragraph{Two distinct primitives, kept distinct.}
The body uses two separate engagement primitives: the \emph{standalone}
improvement $\Delta$ (baseline $\Delta=0.25$), which enters only through
$\Dtil=\rho\Delta$ on the no-takeover branch, and the \emph{premium wedge}
$m_{1}-m_{0}$ (baseline $m_{1}-m_{0}=0.30-0.10=0.20$), which is its own
primitive. The bargaining reading introduces an engagement value
$\Delta_{\mathrm{eng}}$ that shifts the target bloc's \emph{sale} threat point
and so drives the wedge; conceptually $\Delta_{\mathrm{eng}}$ is the engagement
improvement that survives into a sale, and it need not numerically equal the
body's standalone $\Delta$. We keep the two symbols distinct and never conflate
them.

\subsection{Environment and the bargaining problem}
\label{app:bg-env}

Condition on the takeover reaching the bargaining table, i.e.\ on the bidder
finding it worthwhile to bid (the bid rule
$\Pi_{B}=\bar S-P+\xi-\bar m-K\ge 0$ is taken \emph{as given} from the body; we
do \emph{not} re-derive it). Let $a\in\{0,1\}$ index whether engagement has
\emph{succeeded} ($a=1$) or not ($a=0$) by the time control is contested. Two
parties bargain over the gains from transferring control: the target's incumbent
control bloc (Nash weight $\theta\in(0,1)$) and the bidder (weight $1-\theta$).
The dispersed minority float free-rides at the negotiated price and does not sit
at the bargaining table.

\paragraph{The surplus, and why the price $P$ does not enter it.}
If a deal is struck, the value created over and above the parties' disagreement
values is the takeover pie
\begin{equation}
  \label{eq:bg-pie}
  G(a,\xi)\;=\;\bar S+\xi-K-d(a),
\end{equation}
where $\bar S+\xi$ is the gross valuation the bidder realizes, $K$ is the body's
bid cost, and $d(a)$ is the target bloc's disagreement payoff. The cost $K$
enters \eqref{eq:bg-pie} exactly once and is not re-added later. The body's
bidder surplus $\Pi_{B}=\bar S-P+\xi-\bar m-K$ contains the market price $P$,
whereas \eqref{eq:bg-pie} contains \emph{no} $P$; the two are consistent. In the
body the bidder pays $P$ to the \emph{market} --- a transfer to the dispersed
minority float, who free-ride at the negotiated price --- and pays the premium to
the \emph{bloc}. The pie \eqref{eq:bg-pie} is split only between the bidder and
the \emph{bloc}; $P$ is a transfer to a third party (the float) and is orthogonal
to the bidder--bloc split. Formally, total bidder outlay
$=P\cdot(\text{float share})+\Pi_{\text{bloc}}\cdot(\text{bloc share})$; the Nash
solution determines only $\Pi_{\text{bloc}}$, and $P$ enters the body's $\Pi_{B}$
as the separate float-transfer leg. Hence $P$ never enters $G$, and there is no
inconsistency between the body's $\Pi_{B}$ and the appendix's bidder net gain
$\Pi_{A}$.

\paragraph{The disagreement payoff: stated as a condition.}
We use
\begin{equation}
  \label{eq:bg-disagree}
  d(a)\;=\;w+a\,\lambda\,\Dtil_{\mathrm{eng}},
  \qquad \Dtil_{\mathrm{eng}}:=\rho\,\Delta_{\mathrm{eng}},
  \qquad \lambda\in(0,1],
\end{equation}
where $w$ is the target's standalone continuation value absent engagement,
$\Dtil_{\mathrm{eng}}=\rho\Delta_{\mathrm{eng}}$ is the expected engagement
improvement reaching the sale, and $\lambda$ is the fraction of that improvement
the bloc can appropriate at the disagreement node. Equation~\eqref{eq:bg-disagree}
is the single load-bearing modeling object. \textbf{We adopt it as an explicit
structural condition, not as a derived result} (Condition~\ref{cond:bg}); the
free-rider heuristic of \S\ref{app:bg-threat} motivates its sign but does not
constitute a proof of $\lambda>0$ from deeper primitives.

\begin{condition}[Appropriable engagement threat-point shift]
\label{cond:bg}
At the bargaining table the target bloc's disagreement (no-sale-to-this-bidder)
payoff takes the form \eqref{eq:bg-disagree} with a strictly positive
appropriability coefficient $\lambda\in(0,1]$ and a strictly positive expected
engagement improvement $\Dtil_{\mathrm{eng}}=\rho\Delta_{\mathrm{eng}}>0$. That
is, a fraction $\lambda$ of the engagement improvement that survives to a sale is
captured by the bloc as an outside option, so the engagement-induced shift in the
bloc's threat point is $d(1)-d(0)=\lambda\,\Dtil_{\mathrm{eng}}>0$.
\end{condition}

\paragraph{The Nash split and the realized minority premium.}
The bidder's disagreement value is its no-deal payoff $0$, so the bidder's net
gain is $(1-\theta)G(a,\xi)$ and the bloc captures $d(a)+\theta G(a,\xi)$. The
control premium paid is the bloc's bargained rent net of the common baseline $w$:
\begin{equation}
  \label{eq:bg-rent}
  \mR(a,\xi)\;=\;\bigl[d(a)+\theta\,G(a,\xi)\bigr]-w
  \;=\;a\,\lambda\,\Dtil_{\mathrm{eng}}
       +\theta\bigl(\bar S+\xi-K-w-a\,\lambda\,\Dtil_{\mathrm{eng}}\bigr).
\end{equation}
Collecting terms,
\begin{equation}
  \label{eq:bg-rent2}
  \mR(a,\xi)\;=\;\theta\,(\bar S-K-w)\;+\;\theta\,\xi\;+\;a\,(1-\theta)\,\lambda\,\Dtil_{\mathrm{eng}} .
\end{equation}
The economic reading of the last term: because the bloc can walk to its
engagement-improved standalone $d(1)=w+\lambda\Dtil_{\mathrm{eng}}$ (under
Condition~\ref{cond:bg}), it keeps that improvement \emph{net of} the share
$\theta$ already counted in the gross split, i.e.\ a fraction $1-\theta$ of
$\lambda\Dtil_{\mathrm{eng}}$.

\begin{lemma}[Closed-form Nash rent and participation]
\label{lem:bg-split}
\emph{Given} Condition~\ref{cond:bg}, on the consummation set
$\{G(a,\xi)\ge0\}=\{\xi\ge\xi^{\ast}(a)\}$ with $\xi^{\ast}(a)=d(a)+K-\bar S$, the
generalized-Nash split gives the bloc the rent \eqref{eq:bg-rent} and leaves the
bidder $\Pi_{A}=(1-\theta)G(a,\xi)\ge0$. The bloc's individual rationality
($\mR(a,\xi)\ge d(a)-w$, with equality only at $\xi=\xi^{\ast}(a)$) and the
bidder's incentive compatibility (no rent below $d(a)$ accepted; for
$\theta\in(0,1)$ the Nash split is the unique protocol outcome) both hold, so
$\{\xi\ge\xi^{\ast}(a)\}$ is the equilibrium agreement region.
\end{lemma}

\begin{proof}
The maximand $N(x)=x^{\theta}(G-x)^{1-\theta}$ on $x\in[0,G]$ is strictly
log-concave with first-order condition $\theta/x=(1-\theta)/(G-x)$, i.e.\
$x^{\ast}=\theta G$; adding back $d(a)$ gives bloc payoff $d(a)+\theta G$ and the
bidder $(1-\theta)G$. On $\{G\ge0\}$, $\theta G\ge0$ (zero only at $G=0$) and
$(1-\theta)G\ge0$, giving IR and participation. A rent below $d(a)$ violates the
bloc's IR; offering exactly $d(a)$ is the $\theta=0$ corner, excluded by
$\theta>0$, so the unique split is interior. The boundary $\xi^{\ast}(a)$ solves
$G=0$. Numerically: the argmax matches $\theta G$ to $6.0\times10^{-10}$ (Part~A)
and $\texttt{bad\_event}=\texttt{bad\_IR}=\texttt{bad\_IC}=0$ over the simulated
draws per $a$ (Part~B) in \texttt{d4\_bargaining\_check.py}.
\end{proof}

\noindent\emph{Status.} Lemma~\ref{lem:bg-split} is fully proved: it is the
standard generalized-Nash solution under the \emph{stated} disagreement payoff,
and its content is conditional on Condition~\ref{cond:bg} only through the value
of $d(a)$, not through any unproved sign claim.

\subsection{Motivating the threat-point condition: a free-rider tender heuristic}
\label{app:bg-threat}

We now give the economic \emph{motivation} for Condition~\ref{cond:bg}. This
subsection is a heuristic that makes the sign of the threat-point shift plausible;
\textbf{it is not a derivation of $\lambda>0$ from primitives}, and we do not
present it as one (Remark~\ref{rem:bg-threat-honest} states precisely why).

\paragraph{Why ``no sale'' is not a literal status quo.}
``No sale to this bidder'' is not a literal status quo $w$ for the engaging bloc;
it is a continuation value, which is the bloc's outside option in the sense of
\citet{BinmoreRubinsteinWolinsky1986}. Suppose bargaining with the incumbent
bidder breaks down. At the disagreement node the target is a free-standing firm
whose fundamentals have been improved by engagement to $w+a\,\Delta^{\!*}$, where
$\Delta^{\!*}$ is the realized fundamental improvement at the disagreement node. A
competitive fringe of potential acquirers may then mount a takeover. By the
free-rider logic of \citet{GrossmanHart1980}, atomistic shareholders tender only
at a price equal to their post-takeover per-share value, so a value-improving
fringe raider captures essentially none of the improvement; the dispersed float
retains it. The control bloc, holding a pivotal block, is \emph{not} a
price-taker in the tender and retains its pro-rata claim on the improved
fundamentals. This makes it \emph{plausible} that the bloc's continuation value
per share is
\begin{equation}
  \label{eq:bg-threat-motiv}
  d(a)\;\approx\;w\;+\;a\,\Lambda\,\Delta^{\!*},
\end{equation}
with $\Lambda\in(0,1]$ the fraction of the fundamental improvement that survives
to the bloc at the disagreement node. Writing
$\lambda\,\Dtil_{\mathrm{eng}}:=\Lambda\,\Delta^{\!*}$ identifies the coefficient
of Condition~\ref{cond:bg}. The heuristic delivers the \emph{sign} ($\Lambda>0$
whenever the bloc is pivotal and the float will not tender below value); it does
\emph{not} pin $\Lambda$, which depends on block size, tender mechanics, and the
fringe's composition.

\begin{remark}[Honest status: this is motivation, not a closed derivation]
\label{rem:bg-threat-honest}
We emphasise, against the temptation to over-claim, that
\eqref{eq:bg-threat-motiv} is a heuristic and Condition~\ref{cond:bg} is an
\emph{assumption}. A genuine derivation would solve a complete tender-game
equilibrium and produce $\Lambda$ (equivalently $\lambda$) as a function of
primitives (block size, reservation values, the fringe's bidding protocol); that
game is orthogonal to this paper's disclosure mechanism and we do not solve it.
Five earlier attempts to close this step failed for exactly this reason: every
candidate ``derivation'' ended in a free parameter ($\lambda$, or equivalently
the bloc's appropriable share) that is the assumption in another guise. We
therefore treat $\lambda\in(0,1]$ as a reduced-form pivotality parameter and
carry it explicitly everywhere below. The substantive interpretive content ---
engagement enters the bloc's threat point with a strictly positive, appropriable
coefficient, and a recognisable free-rider mechanism makes that sign plausible
--- is real; the claim that this sign is \emph{proved} from primitives is one we
do not make. The cite for the free-rider mechanism is \citet{GrossmanHart1980};
the outside-option reading of the threat point is from
\citet{BinmoreRubinsteinWolinsky1986}.
\end{remark}

\subsection{The wedge is \texorpdfstring{$\xi$}{xi}-free; the levels are not}
\label{app:bg-wedge}

The engagement \emph{wedge} --- the premium object the body's $\kappa$-comparative
statics use --- is the within-state ($\text{common-}\xi$) difference of
\eqref{eq:bg-rent2}:
\begin{equation}
  \label{eq:bg-wedge}
  \mR(1,\xi)-\mR(0,\xi)\;=\;(1-\theta)\,\lambda\,\Dtil_{\mathrm{eng}}
  \;=\;(1-\theta)\,\lambda\,\rho\,\Delta_{\mathrm{eng}} .
\end{equation}
The $\theta\,\xi$ term in \eqref{eq:bg-rent2} is common to both states and
cancels; the bid cost $K$ and baseline $w$ cancel for the same reason. Hence the
wedge is exactly $\xi$-free, $K$-free, and $w$-free. The cancellation is local to
the \emph{difference}: the \emph{level} $\mR(a,\xi)$ is manifestly $\xi$-dependent
through $\theta\,\xi$ (analysed in \S\ref{app:bg-levels}). \emph{Status: proved}
(an algebraic difference of \eqref{eq:bg-rent2}, conditional on
Condition~\ref{cond:bg} only through the value of $\lambda\Dtil_{\mathrm{eng}}$).

\subsection{Reduced form, folding, and the \texorpdfstring{$\rho$}{rho} non-identification}
\label{app:bg-reduced}

The body posits primitives $(m_{0},m_{1})$ and sets
$\mtil=m_{0}+\rho(m_{1}-m_{0})$, so $\mtil-m_{0}=\rho(m_{1}-m_{0})$. The
bargaining reading pins the \emph{wedge} via \eqref{eq:bg-wedge}; the body's
single-$\rho$ fold then reads off $\mtil$ with the \emph{same} $(m_{0},m_{1})$:
\begin{equation}
  \label{eq:bg-identify}
  \boxed{\;m_{1}-m_{0}\;:=\;\mR(1)-\mR(0)\;=\;(1-\theta)\,\lambda\,\rho\,\Delta_{\mathrm{eng}},
  \qquad
  \mtil-m_{0}\;=\;\rho\,(m_{1}-m_{0})\;=\;\rho^{2}(1-\theta)\,\lambda\,\Delta_{\mathrm{eng}}\;}
\end{equation}
The $\rho^{2}$ is justified by an explicit event tree
(Lemma~\ref{lem:bg-tree}, \S\ref{app:bg-tree}): one $\rho$ is the success
probability of the engagement \emph{campaign} acting on the engagement
improvement (inside $\Dtil_{\mathrm{eng}}=\rho\Delta_{\mathrm{eng}}$), the second
is the body's own $\rho$ in the fold (the premium-state arrival acting on the
wedge); the two attach to \emph{distinct independent} Bernoulli draws, so
$\rho^{2}=\E[B_{1}B_{2}]\ne\E[B_{1}^{2}]=\rho$. We \emph{never} divide the
bargained wedge by $\rho$.

\begin{proposition}[Reduced form is the constant-$\theta$ special case]
\label{prop:bg-reduced}
Fix $\theta\in(0,1)$ constant and $\lambda\in(0,1]$, and assume
Condition~\ref{cond:bg}. Then the bargained premia
\eqref{eq:bg-wedge}--\eqref{eq:bg-identify} reproduce the body's reduced form:
there exist $(m_{0},m_{1})$ with
$m_{1}-m_{0}=(1-\theta)\lambda\rho\Delta_{\mathrm{eng}}$ and
$\mtil=m_{0}+\rho(m_{1}-m_{0})$. Conversely, any body calibration
$(m_{0},m_{1},\rho)$ with $m_{1}>m_{0}$ is realized by a one-parameter family of
primitives,
\begin{equation}
  \label{eq:bg-locus}
  \Bigl\{(\theta,\lambda,\Delta_{\mathrm{eng}}):\;
    (1-\theta)\,\lambda\,\rho\,\Delta_{\mathrm{eng}}=m_{1}-m_{0},\ \theta\in(0,1),\ \lambda\in(0,1]\Bigr\},
\end{equation}
e.g.\ $\theta=1-(m_{1}-m_{0})/(\lambda\rho\Delta_{\mathrm{eng}})$ for any
$\Delta_{\mathrm{eng}}>(m_{1}-m_{0})/(\lambda\rho)$. The body verbatim is the case
$\lambda=1$ with state-independent $\theta$.
\end{proposition}

\begin{proof}
Constancy of $\theta$ in $a$ makes the wedge
$(1-\theta)\lambda\rho\Delta_{\mathrm{eng}}$ a single number $m_{1}-m_{0}$; the
body's fold defines $\mtil$. For the converse, given $(m_{1}-m_{0},\rho)$ and any
$\lambda\in(0,1]$, pick any $\Delta_{\mathrm{eng}}>(m_{1}-m_{0})/(\lambda\rho)$
and set $\theta=1-(m_{1}-m_{0})/(\lambda\rho\Delta_{\mathrm{eng}})\in(0,1)$; then
$(1-\theta)\lambda\rho\Delta_{\mathrm{eng}}=m_{1}-m_{0}$.
\end{proof}

\noindent\emph{Status: proved} (nesting plus converse; the only structural input
is Condition~\ref{cond:bg}, which supplies the wedge form).

\begin{remark}[Identification: only the wedge \emph{product} is pinned]
\label{rem:bg-ident}
Proposition~\ref{prop:bg-reduced} is also an honest \emph{non}-identification
statement. The body's calibration pins the product
$(1-\theta)\lambda\rho\Delta_{\mathrm{eng}}=m_{1}-m_{0}$ but not its factors. Two
consequences: (i) the bargaining weight $\theta$ is \emph{set-identified}, not
point-identified, given the wedge and $\rho$ (it traces the locus
\eqref{eq:bg-locus}); and (ii) the structural comparative static in $\theta$
(Remark~\ref{rem:bg-thetadomain}) is therefore conditional on an external pin for
$\theta$. An identifying restriction --- e.g.\ an observable target-resistance
moment mapping to $\theta$, such as the empirical pass rate of management
defenses or an average premium-to-pre-bid ratio in a comparable sample --- would
restore point identification; absent such a moment, only the product is
recoverable. $\rho$ is likewise not separately identified (it enters observables
only through $\Dtil_{\mathrm{eng}}=\rho\Delta_{\mathrm{eng}}$ and the wedge), so
the body's $\rho=0.9$ is a normalization on the locus, not an over-identifying
restriction.
\end{remark}

\paragraph{Reproducing the body's numbers ($\lambda$-explicit reconciliations,
one fold).}
The body's baseline has $m_{0}=0.10$, $m_{1}=0.30$, $\rho=0.9$ --- hence wedge
$m_{1}-m_{0}=0.20$ and, by the body's fold, $\mtil-m_{0}=\rho(m_{1}-m_{0})=0.18$
(so $\mtil=0.28$). With $\Delta_{\mathrm{eng}}=(m_{1}-m_{0})/[(1-\theta)\lambda
\rho]$:
\begin{align}
  \textbf{(R1)}\ \theta=\tfrac12,\ \lambda=1 \ &\Longrightarrow\
    \Delta_{\mathrm{eng}}=\frac{0.20}{(0.5)(1)(0.9)}=\tfrac49\approx0.444,
    \label{eq:bg-R1}\\
  \textbf{(R2)}\ \Delta_{\mathrm{eng}}=\Delta=0.25,\ \lambda=1 \ &\Longrightarrow\
    \theta=1-\frac{0.20}{(1)(0.9)(0.25)}=\tfrac19\approx0.111,
    \label{eq:bg-R2}\\
  \textbf{(R3)}\ \theta=\tfrac12,\ \lambda=\tfrac12 \ &\Longrightarrow\
    \Delta_{\mathrm{eng}}=\frac{0.20}{(0.5)(0.5)(0.9)}=\tfrac89\approx0.889.
    \label{eq:bg-R3}
\end{align}
All three give $(1-\theta)\lambda\rho\Delta_{\mathrm{eng}}=0.20=m_{1}-m_{0}$ and,
through the body's single-$\rho$ fold, $\mtil-m_{0}=\rho\cdot0.20=0.18$. (R1) is
the symmetric-power, full-appropriability benchmark; (R2) is the reading in which
the engagement value is identical on both continuations, predicting a low target
weight $\theta=1/9$; (R3) shows that halving appropriability ($\lambda=1/2$)
\emph{doubles} the required $\Delta_{\mathrm{eng}}$ (from $4/9$ to $8/9$),
reproducing the same wedge. All three are observationally identical at the body's
$0.20$ --- the content of Remark~\ref{rem:bg-ident}: they differ by up to a factor
$\approx4.5$ in $\theta$ and a factor $2$ in $\lambda$ yet generate the same wedge.
Verified this session (Part~D0 and a stdlib cross-check): $(1-\theta)\lambda\rho
\Delta_{\mathrm{eng}}=0.200000$ for (R1),(R2),(R3) and $\rho\cdot0.20=0.18$.

\subsection{The two \texorpdfstring{$\rho$}{rho}'s: an explicit decision/event tree}
\label{app:bg-tree}

The folded object $\mtil-m_{0}=\rho^{2}(1-\theta)\lambda\Delta_{\mathrm{eng}}$
squares $\rho$. The square must attach to \emph{two distinct} Bernoulli events,
or it is the double-deflation bug. We write the tree explicitly. Three sequential,
logically distinct random events govern the activism payoff:
\begin{enumerate}[label=(E\arabic*),leftmargin=2.6em]
\item \textbf{Engagement success} (the $\rho$ inside $\Dtil_{\mathrm{eng}}$).
  Conditional on the blockholder choosing $a=1$, the campaign succeeds with
  probability $\rho$ and \emph{realizes} the fundamental improvement
  $\Delta_{\mathrm{eng}}$, so the success-contingent improvement entering the
  threat point is $\Dtil_{\mathrm{eng}}=\rho\Delta_{\mathrm{eng}}$. Draw
  $B_{1}\sim\mathrm{Bern}(\rho)$.
\item \textbf{Premium-state arrival / fold from $m_{0}$ to $\mtil$} (the body's
  $\rho$ in $\mtil=m_{0}+\rho(m_{1}-m_{0})$). This is the probability that the
  premium-relevant takeover state materializes, pricing the move from $m_{0}$ to
  $m_{1}$. Draw $B_{2}\sim\mathrm{Bern}(\rho)$, independent of $B_{1}$ in the
  body's timing (campaign success and premium-state arrival are separate
  resolutions in the body's tree).
\item \textbf{A bid arrives}, indicator $\1\{\text{bid}\}$ with probability $p$.
  This is the order-flow/disclosure-driven bidding event of the body; it is
  \emph{not} a $\rho$ event and multiplies the level in $\Delta^{\min}$
  separately. Draw $B_{3}\sim\mathrm{Bern}(p)$.
\end{enumerate}

\begin{lemma}[The two $\rho$'s attach to distinct draws]
\label{lem:bg-tree}
With $B_{1}\perp B_{2}$ as in \textup{(E1)--(E2)}, the ex-ante engagement
contribution to the premium is
\[
  \E\bigl[(\mR(B_{1},\cdot)-m_{0})\cdot\1\{B_{2}\}\bigr]
   \;=\;\underbrace{\PP(B_{2}{=}1)}_{=\rho}\cdot
        \underbrace{(1-\theta)\lambda\,\E[B_{1}]\,\Delta_{\mathrm{eng}}}_{=(1-\theta)\lambda\rho\Delta_{\mathrm{eng}}}
   \;=\;\rho^{2}(1-\theta)\lambda\,\Delta_{\mathrm{eng}},
\]
so the two factors of $\rho$ are $\PP(B_{2}{=}1)$ and $\E[B_{1}]$ --- two distinct
Bernoulli draws, not a repeated deflation of one event.
\end{lemma}
\begin{proof}
The wedge realized when the campaign succeeds is
$m_{1}-m_{0}=(1-\theta)\lambda\Dtil_{\mathrm{eng}}$, but $\Dtil_{\mathrm{eng}}$
already integrates out $B_{1}$:
$\Dtil_{\mathrm{eng}}=\E[B_{1}]\Delta_{\mathrm{eng}}=\rho\Delta_{\mathrm{eng}}$
(E1). Multiplying by the fold probability $\PP(B_{2}{=}1)=\rho$ (E2, independent
of E1) gives the displayed product. Each $\rho$ is the mean of a separate,
independent Bernoulli, so the square is the product of two distinct expectations,
not $\E[B^{2}]$ for a single $B$.
\end{proof}

\noindent\emph{Status: proved} (an expectation under the stated independent
event tree; no unproved sign claim enters).

\begin{remark}[Why this is not double-deflation]
\label{rem:bg-nodouble}
The identity $\E[B_{1}B_{2}]=\rho^{2}\ne\E[B_{1}^{2}]=\rho$ is the formal test.
Double deflation would multiply the \emph{same} event's probability twice,
yielding $\E[B_{1}^{2}]=\rho$ (a single $\rho$). The model multiplies two
\emph{independent} events, yielding $\rho^{2}$. The economic distinction: E1 is
``did the engagement campaign work?''; E2 is ``did the premium-relevant takeover
state arrive?''. If a future calibration distinguished their probabilities
($\rho_{1}\ne\rho_{2}$) the fold would read
$\rho_{2}(1-\theta)\lambda\rho_{1}\Delta_{\mathrm{eng}}$, making the two draws
visibly separate; setting $\rho_{1}=\rho_{2}=\rho$ is a calibration choice, not a
conflation.
\end{remark}

\subsection{Assumption A3 as a conditional theorem}
\label{app:bg-A3}

\begin{theorem}[A3 from bargaining, conditional on Condition~\ref{cond:bg}]
\label{thm:bg-A3}
\emph{Assume Condition~\ref{cond:bg}.} Suppose further that the bloc is not a
monopolist over the surplus ($\theta<1$). Then the bargained premia satisfy
\[
  \mtil-m_{0}\;=\;\rho^{2}(1-\theta)\,\lambda\,\Delta_{\mathrm{eng}}\;>\;0,
  \qquad\text{equivalently}\qquad \mtil>m_{0}.
\]
Thus standing Assumption~\textup{(A3)} holds as a \emph{consequence} of bargaining
\emph{given} the structural Condition~\ref{cond:bg}. The condition is sharp: A3
fails iff $\theta=1$ (the bloc captures the entire pie), $\lambda=0$ (engagement
is unappropriable at the threat point, i.e.\ Condition~\ref{cond:bg} is violated),
or $\rho\Delta_{\mathrm{eng}}=0$ (engagement is worthless).
\end{theorem}

\begin{proof}
By \eqref{eq:bg-identify},
$\mtil-m_{0}=\rho^{2}(1-\theta)\lambda\Delta_{\mathrm{eng}}$. Under
Condition~\ref{cond:bg} ($\lambda>0$, $\rho\Delta_{\mathrm{eng}}>0$) and
$\theta<1$ each factor is strictly positive, so the product is strictly positive
and $\mtil>m_{0}$. Sharpness: if $\theta=1$ then $1-\theta=0$ and the wedge
\eqref{eq:bg-wedge} vanishes; if $\lambda=0$ or $\rho\Delta_{\mathrm{eng}}=0$ then
$\lambda\Dtil_{\mathrm{eng}}=0$ and again the wedge vanishes. The wedge form
\eqref{eq:bg-wedge} is verified to machine precision (Part~C, err
$2.1\times10^{-16}$).
\end{proof}

\begin{remark}[What ``conditional theorem'' rests on; sign vs.\ magnitude in $\lambda$]
\label{rem:bg-A3honest}
Theorem~\ref{thm:bg-A3} is a theorem \emph{conditional on}
Condition~\ref{cond:bg}; it is \emph{not} an unconditional theorem and we do not
present it as one. The economic content of A3 has moved from the reduced-form
assertion ``$m_{1}>m_{0}$'' to the structural condition
``$d(1)-d(0)=\lambda\Dtil_{\mathrm{eng}}>0$,'' which the free-rider heuristic of
\S\ref{app:bg-threat} makes plausible but does not prove from primitives. The
\emph{sign} of the wedge holds for any $\lambda>0$; the \emph{magnitude} scales
with $\lambda$ (see (R3)), so the display
$\mtil-m_{0}=\rho^{2}(1-\theta)\lambda\Delta_{\mathrm{eng}}$ carries $\lambda$ and
any quantitative match holds $\lambda$ fixed. We do not claim to have eliminated
all primitives, only to have replaced an opaque premium ordering with a
transparent, motivated threat-point condition.
\end{remark}

\paragraph{Illustrative margin.}
At the headline reconciliation (R1), $(\theta,\lambda,\rho,\Delta_{\mathrm{eng}})=
(\tfrac12,1,0.9,\tfrac49)$ gives $m_{1}-m_{0}=0.20$ and
$\mtil-m_{0}=\rho^{2}(1-\theta)\lambda\Delta_{\mathrm{eng}}=0.18>0$, so A3 holds
with margin $0.18$ --- the body's value.

\begin{remark}[Bargaining power enters the premium --- conditional on a pin]
\label{rem:bg-thetadomain}
On $\theta\in(0,1)$ the bloc captures share $\theta$ of the gross control rent
while the minority benefit of engagement is $(1-\theta)\lambda\rho
\Delta_{\mathrm{eng}}$. \emph{If $\theta$ were independently identified}, this
would deliver a structural comparative static absent from the reduced form:
takeover premia and the minority gain move oppositely in $\theta$, since a more
powerful bloc extracts more of the gross surplus as private rent
($\theta(\bar S+\xi-K-w)$ rises) and concedes less of the threat-shift as minority
value ($(1-\theta)\lambda\Dtil_{\mathrm{eng}}$ falls) --- the control-rent--split
channel of \citet{GrossmanHart1980}. Because $\theta$ is only \emph{set}-identified
(Remark~\ref{rem:bg-ident}), we state this as a conditional structural prediction,
\emph{not} as a calibrated result, pending an external moment for $\theta$. The
toehold/auction extension, in which the engagement stake confers an auction
advantage, is in the spirit of \citet{BulowHuangKlemperer1999}.
\end{remark}

\subsection{The bargaining content of the coincidence lemma}
\label{app:bg-coincide}

\begin{lemma}[Bid-event coincidence]
\label{lem:bg-coincide}
For any constant $c$, the event ``a deal is struck'' with $d(a)=c$ coincides with
$\{\xi\ge c+K-\bar S\}$; in particular it does not depend on the split $\theta$.
Setting $c=d(a)$ reproduces the body's bid event $\{\Pi_{B}\ge0\}$ at the relevant
rent and price.
\end{lemma}
\begin{proof}
$\{\bar S+\xi-K-c\ge0\}=\{\xi\ge c+K-\bar S\}$ is an algebraic identity for any
constant $c$. Verified for arbitrary $c\in\{-0.7,0.123,0.9,2.5\}$ with zero
mismatches over the simulated draws each (spot-check).
\end{proof}

\begin{remark}[Honesty: Lemma~\ref{lem:bg-coincide} is a tautology]
\label{rem:bg-honesty}
Because the coincidence holds for an \emph{arbitrary} constant $c$, it carries no
bargaining content: $\theta$ does not appear, since the \emph{event} of a deal
depends only on whether the pie is positive, not on how it is split. It is a
consistency check that the bargained rent slots into the existing bid rule without
contradiction, \emph{not} a derivation of the body's bid probability $p$. All
bargaining content is in \eqref{eq:bg-wedge}--\eqref{eq:bg-identify} and
Theorem~\ref{thm:bg-A3}.
\end{remark}

\subsection{The level reading and the shape of \texorpdfstring{$\Delta^{\min}$}{Dmin}}
\label{app:bg-levels}

The body's headline object
$\Delta^{\min}(\kappa)=\E[\mR(a)\,\1\{\text{bid}\}]$ uses the \emph{level}
$\mR(a)$, not just the wedge; its decomposition (draft eq.~\eqref{eq:decomp})
contains a level term $m_{0}\PP(\text{bid})$ and an activism term carrying the
wedge $(\mtil-m_{0})$. We must show the $\xi$-dependence of the level in
\eqref{eq:bg-rent2} does not contaminate the \emph{shape} of $\Delta^{\min}$ ---
the single peak at $\kappa^{\ast}=0.59$ and the endpoint symmetry
$\Delta^{\min}(0)=\Delta^{\min}(1)$.

\begin{proposition}[Level enters $\Delta^{\min}$ as $\kappa$-independent
coefficients times analysed bid statistics]
\label{prop:bg-levelshape}
Write the bargaining level as
$\mR(a,\xi)=\theta(\bar S-K-w)+\theta\xi+(1-\theta)\lambda a\,\Dtil_{\mathrm{eng}}$
from \eqref{eq:bg-rent2}. Then
\begin{equation}
  \label{eq:bg-Dmin-decomp}
  \Delta^{\min}(\kappa)
   = \theta(\bar S-K-w)\,\underbrace{\E[\1\{\mathrm{bid}\}]}_{=\,p(\kappa)}
   \;+\; \theta\,\underbrace{\E[\xi\,\1\{\mathrm{bid}\}]}_{=\,\sigma_\xi\,\phi(t(\kappa))}
   \;+\; (1-\theta)\lambda\,\Dtil_{\mathrm{eng}}\,\underbrace{\E[a\,\1\{\mathrm{bid}\}]}_{\text{wedge channel}} ,
\end{equation}
where, under the body's bidding rule
$\1\{\mathrm{bid}\}=\1\{\xi\ge\sigma_\xi t(\kappa)\}$ with threshold
$t(\kappa)=(\bar m+K-\bar S+P(\kappa))/\sigma_\xi$ and $\xi\sim N(0,\sigma_\xi^2)$,
the middle term equals $\theta\sigma_\xi\phi(t(\kappa))$. The first and third
terms are exactly the body's level and wedge channels up to the
$\kappa$-\emph{independent} affine rescalings $\theta(\bar S-K-w)$ and
$(1-\theta)\lambda$. Hence $\Delta^{\min}$ is an affine combination, with
$\kappa$-independent coefficients, of three functions of the \emph{same}
equilibrium objects analysed in the body: $p(\kappa)$, $\sigma_\xi\phi(t(\kappa))$,
and the wedge channel $\E[a\1\{\mathrm{bid}\}]$.
\end{proposition}
\begin{proof}
Take $\E[\,\cdot\,\1\{\mathrm{bid}\}]$ of \eqref{eq:bg-rent2} for $a\in\{0,1\}$
and use linearity. The $\xi$ term gives $\theta\E[\xi\1\{\mathrm{bid}\}]$; under
the Gaussian $\xi$ and the threshold rule,
$\E[\xi\1\{\xi\ge\sigma_\xi t\}]=\sigma_\xi\phi(t)$ (standard truncated-normal
first moment, with the missing-mass term vanishing because $\E[\xi]=0$). The
coefficients $\theta(\bar S-K-w)$, $\theta$, $(1-\theta)\lambda\Dtil_{\mathrm{eng}}$
are constants in $\kappa$.
\end{proof}

\noindent\emph{Status: proved} (linearity of the conditional expectation plus the
truncated-normal first moment).

\begin{corollary}[Shape consistency at the calibration]
\label{cor:bg-shape}
The single-peak property and the endpoint symmetry of $\Delta^{\min}$ recorded in
the body (Phase-0: $\kappa^{\ast}=0.59$, peak $0.07756$,
$\Delta^{\min}(0)=\Delta^{\min}(1)$ up to solver noise $\sim7\times10^{-7}$) are
consistent with the bargaining level. Two facts support this. (i) \emph{Endpoint
symmetry:} the middle term is a fixed multiple $\theta\sigma_\xi$ of
$\phi(t(\kappa))$, and at the symmetric calibration $t(\kappa)$ inherits the
endpoint symmetry of $P(\kappa)$ (Phase-0: $P(D{=}0)$ and the bid statistics are
endpoint-symmetric, $\kappa{=}0$ and $\kappa{=}1$ both at the two-point posterior
law $\{0,\bar\pi\}$), so the middle term is itself endpoint-symmetric and does not
break $\Delta^{\min}(0)=\Delta^{\min}(1)$. (ii) \emph{Single peak:} $\theta
\sigma_\xi$ is small relative to the level coefficient $\theta(\bar S-K-w)$ at the
calibration ($\sigma_\xi=0.40$ vs.\ $\bar S-K-w=O(1)$), so the middle term
perturbs but does not relocate the maximizer. This is confirmed by the re-solve
in Numerical Regularity~\ref{nr:bg-statedep} ($\kappa^{\ast}$ unchanged). We label
this a \emph{consistency} statement at the calibration, not a shape-invariance
theorem.
\end{corollary}

\begin{remark}[Reconciling the level and the wedge-only reading]
\label{rem:bg-levels}
An earlier draft adopted a ``wedge-only'' reading and asserted the level was
$\kappa$-irrelevant. Proposition~\ref{prop:bg-levelshape} makes the honest
statement: the level \emph{does} enter $\Delta^{\min}$ (through
$m_{0}\PP(\text{bid})$ and the $\xi$ term), and $\PP(\text{bid})$ \emph{is}
$\kappa$-dependent. What is true is the weaker, sufficient claim that the level
enters only through (a) the already-analysed bid probability $p(\kappa)$ and (b) a
fixed affine image $\sigma_\xi\phi(t(\kappa))$ of the same bid statistic, with
$\kappa$-independent coefficients; it therefore cannot generate a new
$\kappa$-dependence of a form not already present in the body's analysis, and the
numerical check confirms the shape is intact at the calibration.
\end{remark}

\subsection{Does the single peak survive a state-dependent wedge?}
\label{app:bg-statedep}

The constant-$\theta$ reduced form delivers a constant wedge. A natural robustness
question is whether the body's single-peak comparative static of
$\Delta^{\min}(\kappa)$ survives if the wedge varies with the equilibrium state
--- e.g.\ a bargaining weight that tilts toward the bloc when the blockholder's
position is more informative. We perturb the wedge by a state-dependent factor,
\begin{equation}
  \label{eq:bg-statedep-form}
  m_{1}-m_{0}\;\mapsto\;(1-\theta)\lambda\rho\Delta_{\mathrm{eng}}\bigl(1+\lambda_{\mathrm{sd}}\,g(\kappa)\bigr),
  \qquad g(\kappa)=\kappa-\tfrac12,
\end{equation}
with $g$ bounded and $\lambda_{\mathrm{sd}}\in\{-0.4,-0.2,0,+0.2,+0.4\}$, and
re-solve for the equilibrium cutoffs and $\Delta^{\min}(\kappa)$ at each $\kappa$
with the paper's own solver.

\begin{numreg}[Single peak survives a state-dependent wedge --- at the Phase-0
calibration]
\label{nr:bg-statedep}
This is a \emph{numerically verified regularity at the Phase-0 calibration, not a
theorem.} In the Phase-0 baseline calibration, with the \emph{constant} wedge
($\lambda_{\mathrm{sd}}=0$), $\Delta^{\min}(\kappa)$ is single-peaked over
$\kappa\in[0.01,0.99]$ with interior maximizer $\kappa^{\ast}=0.59$ and peak
$\Delta^{\min}(\kappa^{\ast})=0.07756$ on the canonical $99$-point grid; the
amplitude over the highest converged endpoint is $\approx0.0107$ ($\approx13.9\%$
of level). The appendix's own GE re-solve (\texttt{d4\_bargaining\_check.py},
Part~D1, on a coarser solver grid) returns the same single hump:
$\kappa^{\ast}=0.60$, peak $0.077550$, amplitude $0.0106$, with
$\texttt{up}=\texttt{down}=\texttt{True}$. Under the state-dependent wedge
\eqref{eq:bg-statedep-form} (Part~D2), the interior maximizer stays at
$\kappa^{\ast}=0.60$ and the profile remains single-peaked
(\texttt{interior\_hump}$=$\texttt{True}) for \emph{every}
$\lambda_{\mathrm{sd}}\in\{-0.4,-0.2,0,+0.2,+0.4\}$; the peak rises/falls
monotonically with $\lambda_{\mathrm{sd}}$ from $0.072241$
($\lambda_{\mathrm{sd}}=-0.4$) to $0.082860$ ($\lambda_{\mathrm{sd}}=+0.4$), with
amplitudes $0.0097$ to $0.0114$. All Part~D points pass the solver's residual
gate (the script raises on any non-convergence; OVERALL: ALL PASS). The headline
single-hump orientation is \emph{conditional} on the chord/secant curvature
condition (C$^{\ast\ast}$) of Appendix~D.1: because the realized $D=0$ law is the
two-point set $\{0,\bar\pi\}$, the governing curvature is the secant second
difference of $h(\pi)=\pi\,p(\pi)$ across the chord $[0,\bar\pi]$, not a pointwise
curvature; when (C$^{\ast\ast}$) fails the profile inverts to a trough (see
Appendix~D.1). Across-calibration robustness is a separate exercise and
(C$^{\ast\ast}$) remains a stated hypothesis, not a result.
\end{numreg}

\paragraph{Why the perturbation does not move $\kappa^{\ast}$ (analytic backing).}
By Proposition~\ref{prop:bg-levelshape} the wedge enters $\Delta^{\min}$ only
through the channel $(1-\theta)\lambda\Dtil_{\mathrm{eng}}\,\E[a\1\{\mathrm{bid}\}]$.
A state-dependent factor $1+\lambda_{\mathrm{sd}}g(\kappa)$ rescales \emph{that
channel} by $O(\lambda_{\mathrm{sd}})$ without touching the bidding/posterior
objects $p(\kappa)$ and $\sigma_\xi\phi(t(\kappa))$ that locate the peak; for
moderate $\lambda_{\mathrm{sd}}$ the peak height moves but its location does not,
which is exactly what Part~D2 reports. This is a heuristic for the
\emph{location}'s stability, consistent with --- but not a substitute for --- the
direct re-solve.

\paragraph{Endpoint symmetry (a $\kappa\to0,1$ limit statement).}
Because the protocol delivers a fixed, $\kappa$-independent pair $(m_{0},m_{1})$
(Proposition~\ref{prop:bg-reduced}), the bargaining economy is the body's reduced
form at that pair, and endpoint symmetry
$\Delta^{\min}(\kappa\!\to\!0)=\Delta^{\min}(\kappa\!\to\!1)$ is inherited from
the corrected Appendix~D.1 lemma as a $\kappa\to0,1$ \emph{limit} statement (the
realized $D=0$ posterior law is the same two-point set $\{0,\bar\pi\}$,
$\bar\pi=\omega_{Q}/(\omega_{H}+\omega_{Q})$, at both ends). Phase-0:
$\Delta^{\min}(\kappa=10^{-6})=0.07020822$ and
$\Delta^{\min}(\kappa=1-10^{-6})=0.07020753$ agree to $6.9\times10^{-7}$. The
exact endpoints are not numerically solvable equilibria (the solver hits a
grid-edge degenerate region there), so we state symmetry as a limit, not as a
solved endpoint value.

\subsection*{Phase-0 number reconciliation (single tuple)}
\label{app:bg-numrecon}

This appendix cites \emph{one} canonical Phase-0 tuple everywhere (header,
endpoint paragraph, Numerical Regularity~\ref{nr:bg-statedep}):
\[
  (\text{grid},\ \kappa^{\ast},\ \text{peak},\ \text{amplitude})
  \;=\;
  (\text{99-pt over }[0.01,0.99],\ \ 0.59,\ \ 0.07756,\ \ 0.0107).
\]
The appendix's own coarser GE re-solve returns the consistent
$(\kappa^{\ast},\text{peak},\text{amp})=(0.60,0.07755,0.0106)$; the small
$\kappa^{\ast}$ difference ($0.59$ vs.\ $0.60$) is grid resolution, not a
substantive disagreement.

\subsection{What is proved, what is conditional, and what is open}
\label{app:bg-ledger}

\begin{center}
\renewcommand{\arraystretch}{1.25}
\begin{tabular}{@{}p{0.50\textwidth}p{0.44\textwidth}@{}}
\hline
\textbf{Claim} & \textbf{Status} \\
\hline
Threat-point shift $d(1)-d(0)=\lambda\Dtil_{\mathrm{eng}}>0$
  & \textbf{Stated as Condition~\ref{cond:bg}} (an assumption), \emph{motivated}
    by a free-rider tender heuristic (\S\ref{app:bg-threat}); $\lambda>0$ is
    \emph{not} derived from primitives (Remark~\ref{rem:bg-threat-honest}). \\
Closed-form Nash rent $\mR(a,\xi)=d(a)-w+\theta G^{+}$; IR, IC, agreement region
  & \textbf{Proved} from the protocol given Condition~\ref{cond:bg}
    (Lemma~\ref{lem:bg-split}). \\
Wedge $m_{1}-m_{0}=(1-\theta)\lambda\rho\Delta_{\mathrm{eng}}$ is $\xi$-free;
levels are not
  & \textbf{Proved} (\eqref{eq:bg-wedge}); level shape consistency
    (Prop.~\ref{prop:bg-levelshape}, Cor.~\ref{cor:bg-shape}). \\
A3: $\mtil>m_{0}$ iff $\theta<1$ and Condition~\ref{cond:bg}
  & \textbf{Theorem conditional on Condition~\ref{cond:bg}}
    (Theorem~\ref{thm:bg-A3}, Remark~\ref{rem:bg-A3honest}); \emph{not}
    unconditional. \\
Folding $\mtil-m_{0}=\rho^{2}(1-\theta)\lambda\Delta_{\mathrm{eng}}$; $\rho^{2}$ =
two distinct draws
  & \textbf{Proved} (event tree, Lemma~\ref{lem:bg-tree},
    Remark~\ref{rem:bg-nodouble}); no divide-by-$\rho$. \\
Reduced form is the constant-$\theta$, $\lambda{=}1$ special case
  & \textbf{Proved} (nesting + converse, Proposition~\ref{prop:bg-reduced}). \\
Identification: only the wedge \emph{product} is pinned; $\theta$ set-identified
  & \textbf{Honest non-identification} (Remark~\ref{rem:bg-ident}); $\theta$
    comparative static \textbf{downgraded} to an if-identified remark
    (Remark~\ref{rem:bg-thetadomain}). \\
Consummation event coincides with the body's bid rule
  & \textbf{Tautology, no bargaining content} (Lemma~\ref{lem:bg-coincide},
    Remark~\ref{rem:bg-honesty}). \\
Endpoint symmetry $\Delta^{\min}(0)=\Delta^{\min}(1)$
  & \textbf{$\kappa\to0,1$ limit statement} (inherits from Appendix~D.1;
    exact endpoints not solvable). \\
Single interior peak (baseline + state-dependent wedge)
  & \textbf{Numerical Regularity at the Phase-0 calibration}, conditional on
    (C$^{\ast\ast}$); state-dependent half \textbf{verified} (D2, all
    $\lambda_{\mathrm{sd}}$ PASS, Numerical Regularity~\ref{nr:bg-statedep});
    \emph{not} a theorem. \\
Price $P$ excluded from the pie $G$
  & \textbf{Reconciled} (\S\ref{app:bg-env}): $P$ is a float transfer orthogonal
    to the bidder--bloc split. \\
\hline
\end{tabular}
\end{center}

\subsection*{Open (precisely what remains)}
\label{app:bg-open}

We state the open points precisely, so that no reader mistakes the conditional
results for unconditional ones.
\begin{enumerate}[label=(O\arabic*),leftmargin=2.6em]
\item \textbf{The threat-point condition is not derived.} Condition~\ref{cond:bg}
  --- specifically the appropriability coefficient $\lambda\in(0,1]$
  (equivalently $\Lambda$) --- is the load-bearing assumption. The free-rider
  heuristic (\S\ref{app:bg-threat}) makes the \emph{sign} $\lambda>0$ plausible
  but does not derive $\lambda$, or even its strict positivity, from a complete
  tender-game equilibrium. A full derivation requires solving the post-engagement
  tender/holdup game (block size, reservation values, the fringe's bidding
  protocol) --- e.g.\ a \citet{GrossmanHart1980} free-rider auction or a
  private-benefit holdup stage in the spirit of the Burkart--Gromb--Panunzi
  literature --- which is orthogonal to this
  paper's disclosure mechanism. This is the step five attempts did not close, and
  we do not claim to have closed it. Consequently A3 (Theorem~\ref{thm:bg-A3}) is
  a \emph{conditional} theorem.
\item \textbf{$\Delta_{\mathrm{eng}}$ is not identified with the body's $\Delta$.}
  Reading (R2) maintains $\Delta_{\mathrm{eng}}=\Delta=0.25$; this is a
  \emph{maintained reading}, not a result. Absent it, $(\theta,\lambda,
  \Delta_{\mathrm{eng}})$ is pinned only up to the wedge product
  (Remark~\ref{rem:bg-ident}), and $\theta$ is only set-identified.
\item \textbf{The single-peak result is numerical, not analytic.} The interior
  peak and its survival of a state-dependent wedge (Numerical
  Regularity~\ref{nr:bg-statedep}) are verified at the Phase-0 calibration only,
  and the orientation (hump vs.\ trough) is conditional on the chord-curvature
  condition (C$^{\ast\ast}$) of Appendix~D.1. Across-calibration robustness and a
  closed-form sufficient condition for (C$^{\ast\ast}$) remain open.
\item \textbf{Equilibrium uniqueness is numerical.} We take
  Assumption~\textup{(A6)} (cutoff-map contraction) as given, so equilibrium
  uniqueness is numerical at the calibration, not proved.
\end{enumerate}

%==============================================================================
%  SELF-CHECK (Bash / d4_bargaining_check.py, this session):
%   BODY: Delta=0.25 (standalone) ; m1-m0=0.20 (separate primitive) ;
%         mtilde-m0=rho*(m1-m0)=0.18 ; tildeDelta=rho*Delta=0.225.
%   stdlib cross-check (this session):
%     (R1) theta=1/2, lam=1   -> Delta_eng=0.444444, wedge=0.200000      PASS
%     (R2) theta=1/9, lam=1   -> Delta_eng=0.250000, wedge=0.200000      PASS
%     (R3) theta=1/2, lam=1/2 -> Delta_eng=0.888889, wedge=0.200000      PASS
%     rho*wedge = 0.18 = mtilde-m0                                       PASS
%   d4_bargaining_check.py (verified output _d4_check_out.txt, ALL PASS):
%     [A] argmax x*=theta*G : maxabserr 6.0e-10                          PASS
%     [B] consummation=xi>=xi*(a); IR,IC : bad_event=bad_IR=bad_IC=0     PASS
%     [C] A3: m1-m0=(1-theta)*Dtil (err 2.1e-16); mtilde-m0=rho^2(1-theta)Delta PASS
%     [D0] constant-theta special case: mtilde 0.28==0.28               PASS
%     [D1] baseline GE single peak: kstar=0.60 peak 0.07755 amp 0.0106
%          up=True down=True                                            PASS
%     [D2] state-dependent wedge lam in {-0.4,-0.2,0,+0.2,+0.4}:
%          kstar=0.60 for all; peak 0.072241..0.082860; amp 0.0097..0.0114;
%          interior_hump=True for all                                   PASS
%     OVERALL: ALL PASS
%   HONEST RELABEL: the threat-point shift (Condition cond:bg) is STATED, not
%   derived; the free-rider stage is MOTIVATION; A3 (thm:bg-A3) is CONDITIONAL on
%   Condition cond:bg; the single-peak result (nr:bg-statedep) is a NUMERICAL
%   REGULARITY at the Phase-0 calibration, not a theorem. lambda carried
%   explicitly; rho^2 = E[B1 B2] != E[B1^2]=rho (two distinct independent draws);
%   no divide-by-rho anywhere; body standalone Delta and bargaining Delta_eng
%   kept DISTINCT.
%==============================================================================
